\subsection{Isolated shadows and a complete tower of homology refinements}

\begin{statement}{D.26}{Proved}{The CM witnesses survive every abelian-kernel refinement}
Keep $B,Q,N,F$ and $G_j$ from D.25, and let
$\widetilde N_j\triangleleft B_3$ be the inverse image of
$[N,N]N^j$. Every $\widetilde N_j$ is an isolated object of
the gentle GT-shadow groupoid: every shadow with that target
induces an automorphism of $G_j$.

Put $C=N^{\rm ab}\simeq\mathbb Z^{97}$, and let $s,r$ act
on $C$ by conjugation. For $t=sr$ define $c=[t^{12}]\in C$.
The integral maps
\[
 L:C\oplus C\longrightarrow C,\quad
 L(b,a)=(1-s)b+s(1-r)a,\qquad
 \widetilde L(b,a,h)=L(b,a)-hc                         \tag{D.26.1}
\]
have Smith invariant factors $(1^{96},2)$ and $(1^{97})$,
respectively. In particular $\operatorname{coker}L=C_2$
and $\widetilde L$ is onto over $\mathbb Z$.

There is an exact sequence and a specified quotient map
\[
 1\longrightarrow\widehat{\mathbb Z}^{97}
 \longrightarrow\mathcal G:=\varprojlim_j G_j
 \longrightarrow Q\longrightarrow1,\qquad
 \widehat B\twoheadrightarrow\mathcal G
   =\widehat B/\overline{[\widehat N,\widehat N]} .       \tag{D.26.2}
\]
Each of D.25's twenty-four certified automorphisms of $G_6$
lifts to a continuous automorphism of $\mathcal G$ inducing
the identity on $Q$ and sending $t$ to $t^\lambda$, for some
$\lambda\in\widehat{\mathbb Z}^{\times}$ with
$\lambda\equiv1\pmod{72}$. A single cusp gauge supplies a
compatible family of GT shadows at all $\widetilde N_j$.
For each of the twelve extra CM pairs, its member at level six
can be chosen to be one of D.25's nonarithmetical shadows,
with $m=(\lambda-1)/2\equiv0\pmod{36}$.
Thus all twelve extra CM pairs survive the entire family of
abelian-kernel refinements of $N$.

This family does not exhaust the finite quotients in D.22:
there is an explicit finite quotient of $B$ below $G_6$ whose
kernel does not contain $[N,N]$. The existence of a genuine
full profinite GT$_0$ lift of these shadows remains undecided.
\end{statement}
\begin{proof}
First we verify isolation directly from finite shadow equations.
Let $(m,f)$ be a shadow at $\widetilde N_j$, put
$\lambda=2m+1$, and choose a representative $f\in[F,F]$.
The hexagons give a surjection $\phi:B_3\to G_j$ and send
the braid center to its $\lambda$th power, which is trivial
in $G_j$. Hence $\phi$ factors through $B$.
The cusp $t$ has order $12j$ in $G_j$: its image in $Q$ has
order twelve, and the cycle $c$ is primitive in $C$.
The last assertion also follows from the integral cycle
coordinates in the calculation below, one of which is one
and all the others are zero. Thus
$(\widetilde N_j)_{\rm ord}=6j$ and the shadow's multiplier
is a unit modulo twelve.

Evaluate the same word $f$ in the linear matrices $U^2,W^2$
of A.17 modulo twelve. Its reduction modulo four is $I$,
since $U^2,W^2$ generate an abelian four-group modulo four.
Consequently the linear matrices
$A=U^\lambda$ and $D=f^{-1}W^\lambda f$ satisfy the braid
identity with positive sign. Indeed the projective identity
first gives $ADA=\pm DAD$, and modulo four the choice
$\lambda\equiv\pm1$ gives the positive identity for
$U^\lambda,W^\lambda$; an invertible matrix cannot equal
its negative modulo four. The pair $A,D$ generates all of
$\mathrm{SL}_2(\mathbb Z/12)$: its generated subgroup
surjects onto $Q$, and a proper such subgroup would be a
complement to $\{\pm I\}$. That complement cannot exist,
because both lifts $\pm S$ of the order-two element $s\in Q$
have square $-I$.

Write $D=(\begin{smallmatrix}a&b\\c_0&d\end{smallmatrix})$.
The trace, determinant, braid identity and generation now give
\[
 a+d=2,\quad ad-bc_0=1,\quad
 c_0(1+\lambda c_0)=0,\quad c_0\in(\mathbb Z/12)^\times.
\]
Thus $c_0=-\lambda^{-1}$ and the matrix
$V=(\begin{smallmatrix}\lambda&\lambda(1-a)\\0&1\end{smallmatrix})$
conjugates $U,W$ to $A,D$, respectively. This repeats the
elementary matrix argument of A.17 using only the finite
shadow hypotheses. It proves $q\phi=\operatorname{Int}(V)q$.
In particular $\phi(N)\subset C/jC$. Since this subgroup
is abelian of exponent $j$, $\phi$ kills $[N,N]N^j$.
It therefore induces a surjective endomorphism of the finite
group $G_j$, hence an automorphism. Its kernel is precisely
$\widetilde N_j$. This is isolation in
\cite[Definition3.13]{dDolgushevGuynee2024}.

For the integral assertion use the same oriented graph and
spanning tree as in D.25. The ninety-seven fundamental cycles
form a matrix with the identity matrix in the non-tree-edge
coordinates, so projection to those coordinates is an exact
integral inverse on $C$. The columns of $L$ are
$(1-s)C$ and $s(1-r)C$ in this basis. The program
\texttt{d\_cm\_refine.py} computes their integral Smith form
as $(1^{96},2)$. Appending the column $-P(t^{12})$ gives
Smith form $(1^{97})$; an independent Hermite-normal-form
calculation of this augmented matrix returns $I_{97}$.
The program retains the complete integer coordinates and
reproduces the spanning tree. Thus these are integer-image
computations, not deductions from ranks modulo two and three.

We construct the lift, retaining the multiplier's full
residue. Fix one of the twenty-four certified matrix pairs
$A_0,\beta_0$ and choose the fixed integral paths $P_A,P_\beta$
of D.25. Let $E\in C$ be its integral cusp error, the left
side of (D.25.5). With cycle corrections $(b,a)$ the exact
cusp equation for multiplier $\lambda=1+12h$ is
\[
       E+L(b,a)-hc=0.                                  \tag{D.26.3}
\]
To see the exponent, $P(t^{1+12h})=P(t)+hc$, since
$q(t)c=c$. This proves (D.26.3) in every finite quotient,
and by continuity for $h\in\mathbb Z_p$.

At $p=2,3$ begin with D.25's correction modulo $p$ and
$h=0\pmod p$. Suppose coefficients have been constructed
modulo $p^k$; their error in (D.26.3) is divisible by $p^k$.
Surjectivity of $\widetilde L$ supplies a correction
$(\delta b,\delta a,\delta h)$ modulo $p$ equal to the
negative error divided by $p^k$. Add $p^k$ times that
correction. This constructs compatible solutions over
$\mathbb Z_p$. Their induced matrix on $C\otimes\mathbb Z_p$
reduces to D.25's invertible matrix modulo $p$, hence is
invertible. The associated map on the quotient $Q$ is the
identity, so it is an automorphism of the whole extension.
The initial values $h=0\pmod2$ and $h=0\pmod3$ are
preserved. Accordingly the combined multiplier satisfies
$\lambda\equiv1\pmod{72}$, not just modulo twelve.
As a check of the power and sign conventions, the program
carries every one of the twenty-four solutions through
exponents $2,4,8,16$ and $3,9,27,81$, and directly verifies
both general-$m$ hexagons in their faithful group models:
there are384 successful identity checks.

For $p>3$ take the identity automorphism on
$\mathcal G_p=\varprojlim_k G_{p^k}$, with multiplier one,
but choose conjugators with the required images
$A_0,\beta_0$ in $Q$. Such choices exist. For example,
an arbitrary lift $a_0$ of $A_0\in C_Q(r)$ has
$a_0ra_0^{-1}r^{-1}=z\in C\otimes\mathbb Z_p$ and
$(1+r+r^2)z=0$. Multiplication by three is invertible,
and the identity
$(1-r)(2+r)=3-(1+r+r^2)$ shows that
$z=(1-r)((2+r)z/3)$. Adjusting $a_0$ by the negative of
this cycle gives a lift centralizing $r$. The same argument
for $s$ uses $(1-s)z=2z$ when $(1+s)z=0$ and division by
two. These are compatible $p$-adic choices, and both belong
to the preimage of $H=q(F)$.

The extensions $\mathcal G_p$ have the same quotient $Q$,
and their fibre product over $Q$ is precisely $\mathcal G$.
This follows either from the path model or from Chinese
remainders on $C/jC$. The constructed automorphisms and
conjugators therefore assemble to $\alpha\in
\operatorname{Aut}(\mathcal G)$ and $A,\beta\in\overline F$,
with $\alpha(s)=\beta s\beta^{-1}$,
$\alpha(r)=ArA^{-1}$ and $\alpha(u)=u^\lambda$.
Here $\overline F$ denotes the image of $\widehat F$ in
$\mathcal G$. Its abelianization is
$\overline F^{\rm ab}=\widehat{\mathbb Z}x\oplus
\widehat{\mathbb Z}y$, because the kernel
$\overline{[\widehat N,\widehat N]}$ is contained in
$\widehat F'$.

Write $\beta^{\rm ab}=ax+by$, set $k=b-a$, and conjugate
$\alpha$ by $x^k$. With $\beta'=x^k\beta$ and
$f'=s\beta's(\beta')^{-1}$, the same abelianization
calculation as in D.25 gives $f'\in\overline F'$.
The resulting automorphism satisfies
\[
 \alpha'(u)=u^\lambda,\qquad
 \alpha'(v)=(f')^{-1}v^\lambda f',\qquad
 \alpha'(s)=sf'=(f')^{-1}s.                             \tag{D.26.4}
\]
Since $r=uv$, $s=uvu=vuv$, and $m=(\lambda-1)/2$, its
two hexagons are exactly
\[
 u^\lambda(f')^{-1}v^\lambda f'=(f')^{-1}r x^{-m},\qquad
 (f')^{-1}v^\lambda f'u^\lambda=vu y^{-m}f'.             \tag{D.26.5}
\]
Reduction to every $G_j$ gives a GT shadow: the center is
killed, the unit condition holds, $f'$ has commutator image,
and the induced map is an automorphism. The single global
pair makes these reductions compatible.

For precision concerning D.25's finite gauge, choose its
discrete representative of $\beta$ to agree with the
constructed profinite $\beta$ both in $G_6$ and in
$\overline F^{\rm ab}/36$. Such a representative exists
because $F$ is dense in $\overline F$ and this is a finite
joint quotient. Then its integer $b-a$ agrees with $k$
modulo36, the order of $x$ in $G_6$. Thus its D.25 gauge
is exactly the reduction of (D.26.4), with $m=0\pmod{36}$.
The gauge preserves $T=A^{-1}V$ and $W=\beta^{-1}V$.
Consequently all twelve extra CM pairs remain
nonarithmetical at level six in the compatible tower.
This specifies the finite gauges that extend; it does not
make an assertion about an independently chosen gauge.

Finally we locate the remaining finite tests by an actual
quotient. Choose a free basis of the rank97 group $N$ and
map its first two generators to $(12),(23)\in S_3$ and
the others to the identity. Let $J$ be the normal core in
$B$ of the intersection of this map's kernel with
$[N,N]N^6$. Then $J\triangleleft B$ has finite index,
$J\subset[N,N]N^6$, but $[N,N]\not\subset J$, since
the commutator of the first two basis elements has nontrivial
image in $S_3$. The quotient map $B\to B/J$ therefore
cannot factor through any $G_j$: every such factorization
would kill $[N,N]$. This also proves the kernel in
(D.26.2) is nontrivial and that the homology tower is not
cofinal in the system of D.22.

The source search supplies no additional lifting assertion
for these quotients. In particular isolation gives reduction
\emph{homomorphisms}, as stated in
\cite[Remark3.16]{dDolgushevGuynee2024}; the all-finer-quotient
criterion is its Theorem5.2 and Corollary5.4.
The abelian sufficient criterion in
\cite[TheoremB.2]{dDolgushevLeLorenz2024} requires the
pure quotient to be abelian. Here its image already
surjects onto
$H\simeq C_2^2\times\mathrm{SL}_2(\mathbb F_3)$:
the isomorphism sends the unique lift with modulo-four
image in $V_4$ to its modulo-four and modulo-three components.
The latter group's second derived subgroup has order two.
Similarly the dihedral quotients of
\cite[Remark3.2]{dBDHP2026} embed their $F$ images into
products of three dihedral groups and are metabelian.
No such pure quotient can surject onto this $H$, since
surjections preserve the derived-series images.
These explicit maps explain why those sufficient families
do not supply the remaining lift.
\end{proof}
